\section{行政强制}
\begin{enumerate}
  \item 施加暴力去实现管控目标。
  \item 没有听证程序。\yi 没有为当事人创设新的义务。\er 来不及。\al{作出加处罚款决定前应当履行听证程序 \cuo}。
  \item 强制执行。\bs{不履行行政决定}的公民、法人、其他组织，依法\bs{强制其履行义务}的行为。
        \begin{enumerate}
          \item 行政决定（实务说法）= 具体行政行为。先说 \yjt 再等 \yjt 再做。
          \item 拆房子每年必涉及。俩行为。违反城乡规划法走前面，违反土地管理法走后面。
                \begin{enumerate}
                  \item 先说。规划局\zd{责令拆是处罚}。县以上政府农业农村主管部门。\al{镇政府 \cuo}。
                  \item 再做。县政府拆才是执行。申请法院执行。
                  \item 合法建筑的拆除，只能找法院执行。
                \end{enumerate}
          \item 特征。
                \begin{enumerate}
                  \item 依附性。不是上来就直接执行，存在一个基础决定。
                  \item 强制性。动用暴力。
                  \item 目的性。执行。
                \end{enumerate}
          \item 设定。一律交给法律。执行侵害大，需要收紧。\bs{强制执行仅法律}。
          \item 直接强制执行。一步到位。\al{划拨、拍卖、拆房子、税务局强制扣缴 500W 骗税额}。
                \begin{enumerate}
                  \item 主体个别授权。
                        \begin{enumerate}
                          \item 人身类 \yjt 公安、国安。财产划拨 \yjt 税务、海关。\al{银行、银保监会 \cuo}。拆违建 \yjt 县政府。
                          \item 其他行政机关\zd{申请法院强制执行 \yjt 审查 \yjt 准予 \hz 不准予}。
                          \item 拍卖权大家都有，但要山穷水尽（\zd{不履行、不复议、不起诉}），扣的货物抵罚款。\al{抵税款 \cuo}。
                        \end{enumerate}
                  \item 程序。双轨制。有权自己执行、没权申请法院执行。互斥。例外。税务、海关。
                        \al{市税务局可以书面通知银行划扣甲公司存款 \dui}。
                \end{enumerate}

                \begin{kk}{自己执行}
                  \begin{enumerate}
                    \item 只需等满义务履行期，不用等起诉期、3 个月内申请执行。\al{自起诉期届满 \xhx{180} 日提出 \cuo}。
                    \item 步骤。\bs{一催、二辩、三执行}。
                          \begin{enumerate}
                            \item \zd{书面催告}。过程性行为不可诉。
                            \item \zd{陈述申辩}。\al{\xhx{直接}执行 \cuo}。
                            \item \zx{书面决定}强制执行并\zx{送达}当事人。
                          \end{enumerate}
                    \item \zd{比例原则}的限制。
                          \begin{enumerate}
                            \item 方式。\zx{可以达成执行协议/执行和解}。\al{分阶段履行、少交加处罚款和滞纳金}。
                            \item 当事人不履行，恢复强制执行。\al{申请法院强制执行 \cuo{自己就能执行，为啥要找法院}}。
                            \item 时间。原则\zd{不能夜间（22 - 6）、法定节假日（周末、少民节日）}，紧急情况除外。\al{洪水来了}。
                            \item 手段。\zd{不得对居民生活采取停水、停电、停热、停气}。企业可以断。
                            \item 不得采用暴力、胁迫及其他非法方式。
                            \item 特殊标的。房子必须得等满起诉期，不能直接上来就拆。
                          \end{enumerate}
                  \end{enumerate}
                \end{kk}

                \begin{kk}{申请法院执行（非诉执行）}
                  难点，但不是重点，只考过一次。
                  \begin{enumerate}
                    \item 履行期、复议期、起诉期届满（齐头并进算）后 3 个月内 + Opt(催告 10 日)（届满后才催）。
                    \item 裁决案件的权利人也可找法院，但得等行政机关期满。\al{甲乙土地争议，土地局裁决归甲}。
                    \item 管辖。申请人所在地、基层法院。不动产 \yjt 不动产所在地。\al{向该\xhx{公司所在地}法院申请 \cuo}。
                    \item 行政庭、书面审查、1 - 7 日 + 30、特事特办院长批、可以收取合理费用。\al{执行庭 \cuo}。
                  \end{enumerate}
                \end{kk}
          \item 间接强制执行。\bs{代履行、立即代履行、滞纳金、加处罚款}。
                \begin{enumerate}
                  \item \zd{代履行。为了公共利益}，代替履行义务。
                        \begin{enumerate}
                          \item 主体。自己 \hz 第三方。
                          \item 提前 3 日送达决定不适用立即代履行。
                          \item \zd{所需费用由违法者承担}。
                          \item 不能暴力胁迫、可以派人现场监督。
                        \end{enumerate}
                  \item 立即代履行。\al{高速路车掉货、xx 传播了、xx 泄露了}。
                  \item 执行罚。为了执行而来的罚，震慑你。\bdy 罚款。
                        \begin{kk}{滞纳金、加处罚款（罚款领域的滞纳金）}
                          标准要告知；按日 3\%；上限不超本金；最多收 30 日；复议、诉讼期间不计算。
                        \end{kk}
                \end{enumerate}
        \end{enumerate}
  \item 强制措施。为了\bs{制止违法}、避免\bs{证据损毁}，防止\bs{危险的发生}，控制危险的\bs{扩大}，对于当事人的\bs{人身自由 \hz 财务}的一种\bs{暂时性的控制}手段。
        \al{税务稽查局认定甲公司涉嫌转移财产逃税，扣押其相当于应缴税款的商品 \dui{税收保全}}。
        \begin{enumerate}
          \item 为预防而强制，处罚是板上钉钉。\al{疑似病患者，强制隔离}。
          \item 特征。
                \begin{enumerate}
                  \item 行政性。一定是在行政管理过程当中发生的。区别司法强制措施。
                  \item 非惩罚性。处罚是故意制裁你，你一定违法。
                  \item 暂时性。\al{强制约束醉汉至酒醒（目的实现）}。
                  \item 强制性。动用暴力的物理性的限制。
                \end{enumerate}

                \begin{kk}{vs 执行}
                  \begin{enumerate}
                    \item \zd{目的}不同。措施目的是预防，执行目的就是执行。
                    \item \zd{期限}不同。执行是永久的。
                    \item \zd{阶段}不同。早中晚、前中后。措施 \yjt 罚款 \yjt 执行。
                    \item 措施没有基础决定，说做合一，时空同步。预防目的注定了事发突然来不及走程序办手续。
                    \item 一上来就是动用暴力的 \yjt 强制措施。
                  \end{enumerate}
                \end{kk}

                \begin{kk}{vs 处罚。有时间的是处罚，没时间的是措施}
                  \begin{enumerate}
                    \item 暂扣许可证 \yjt 处罚。扣留/扣押许可证 \yjt 强制措施！！爱考。\al{扣留驾驶证的行为为强制措施 \dui}。
                    \item 处罚有期限（永久也是一种期限，\al{如责令关停}）。题目会出现“行为 + 时间”等字眼，\al{如暂扣 6 个月}。
                          强调明确性，都有量和度。
                    \item \zd{目的}不同。制裁性；制止性和预防性。
                    \item \zd{阶段}不同。行政程序终了之时；行政程序前端。
                    \item \zd{表现形式}不同。列举。
                  \end{enumerate}
                \end{kk}
          \item 种类。
                \begin{enumerate}
                  \item 限制人身自由。
                  \item 查封、扣押（小物件）、冻结。
                  \item 其他暂时没名的。\zd{责令改正、责令恢复原状}。不用背太细。
                \end{enumerate}
          \item 设定。排除规章。
                \begin{enumerate}
                  \item 行政法规。2 个不能。不能限制人自由、钱自由（冻结） \yjt 专属于法律。
                  \item 地方性法规。只能查封、扣押。\al{国务院部门\xhx{规章}不能设定扣押 \dui}。
                \end{enumerate}
          \item 实施。\bs{设定、授权 + 行法；行法没有人和钱，地法只有查和扣；集措罚执禁委托}。
                \begin{enumerate}
                  \item 主体。
                        \begin{enumerate}
                          \item 授权。+行政法规。
                          \item 委托。\bs{强制禁止委托}（动用了暴力）。
                                \begin{itemize}
                                  \item \al{工商局\xhx{委托消协}去扣押 \cuo}。
                                  \item \al{市监局\xhx{不得委托}其他机关实施查封和扣押 \dui}。
                                \end{itemize}
                          \item 集中。集中行使了处罚权的综合执法机关可以实施与处罚权有关的强制措施权。
                                \begin{itemize}
                                  \item 处罚权 \yjt 是车，实施强制措施权 \yjt 车钥匙（附随义务）。送车得给车钥匙。
                                  \item 措 \yjt 罚 \yjt 执。只给处罚权，没给相应的措施权（\al{如证据获取}），处罚权会被架空。
                                \end{itemize}
                        \end{enumerate}
                  \item 程序。
                        \begin{enumerate}
                          \item 调取证据。通知到场、说明依据。
                          \item 动用暴力的。原则报批，紧急情况可事后补办，财产 24 小时，人身立即。
                          \item 现场笔录。见证据部分。
                        \end{enumerate}
                  \item 查封、扣押的特别程序。
                        \begin{enumerate}
                          \item 对象。1 个限于（有关系的物品） 3 个不得（与违法活动无关、生活必需品、重复查封）。
                          \item 善意第三人/受害人/无关的 \yjt 登记即可。\al{对甲\xhx{偷开的车}县公安局可以扣押 \cuo}。
                          \item 形式要求。当场交付决定书、清单（一式两份）。简单的不得了但是必考。
                                \al{9/5 扣押，9/8交付扣押决定书和清单。区城管局交付扣押物品决定书和清单的行为违法 \dui{难为出题老师了}}。
                          \item 期限 30 + 30！不收费。\al{最长期限为 90 天 \cuo}。
                          \item 自己保管 \hz 委托第三人保管（性质 = 民事委托保管合同）。
                          \item 后续处置。\yi 有价值充公、没收。\er 没价值依法销毁。\san 无关到期归还原物 \hz 视情况给补偿。
                        \end{enumerate}
                  \item 冻结。先做后说。
                        \begin{enumerate}
                          \item 不能直接操作账户 \yjt 向金融机构交付冻结书 \yjt 银行立即冻结。
                          \item 事后 3 日内再向当事人送达通知书。
                        \end{enumerate}
                \end{enumerate}
        \end{enumerate}
  \item \zd{协调配合机制}。
        \begin{enumerate}
          \item 应当\zd{加强}协调配合。
          \item 建立健全案件\zd{移送}制度。
          \item 加强\zd{证据材料}移交、接收衔接。
          \item 完善案件处理\zd{信息通报}机制。
        \end{enumerate}
\end{enumerate}
